\chapter{Données de marché pour la recherche}\label{ch:rs:market-data-for-research}

Deux chercheurs mesurent la volatilité d'une même action le même jour à partir de rendements à la seconde. L'un utilise les prix des transactions,
l'autre les mid, et la variance du premier est 3.6 fois celle du second. Échantillonnés toutes les cinq minutes, les deux résultats concordent à 2~\% près.
Aucun des deux chercheurs n'a commis d'erreur~; chacun a choisi une représentation du marché, et toute représentation conserve une partie de ce qui
s'est produit et détruit le reste. Ce chapitre porte sur ces choix~: ce qu'un flux de messages enregistre, les horloges selon lesquelles un
flux de transactions est découpé en \omterm{def:rs:market-data-for-research:bar}{barres}, ce que perdent respectivement les prix de transaction, les mid et les \omterm{def:rs:market-data-for-research:bar}{barres}, comment un historique continu se construit à partir
de contrats à terme qui expirent, et quelle quantité de données tout cela représente. Ses données sont une journée de marché simulée par le simulateur de marché du livre,
construite dans ce chapitre, et neuf années de prix de règlement-livraison des contrats à terme sur le pétrole brut.

\section{Ce que contient un flux de messages}

Le livre~1, chapitre~28, a décrit les niveaux de données de marché~: les meilleurs prix (niveau 1), la profondeur par prix (niveau 2)
et le carnet ordre par ordre (niveau 3), avec les horodatages de bourse et de réception, les numéros de séquence et les conditions de transaction. Un entrepôt de recherche
part du plus riche d'entre eux, un flux ordre par ordre, parce que toute vue plus grossière peut être reconstruite à partir de lui et qu'aucune vue plus fine
ne peut être reconstruite à partir d'une vue plus grossière.

Le simulateur du livre, \texttt{firm.tape}, écrit trois types de messages dans le style du flux TotalView-ITCH de Nasdaq (le tableau
ci-dessous). Un ordre est ajouté avec un identifiant, un côté, un prix et une quantité~; il est ensuite annulé en totalité ou en partie, ou exécuté
contre un ordre entrant. Les transactions ne sont pas des faits distincts~: ce sont les exécutions, et le côté de l'agresseur est l'opposé de
celui de l'ordre au repos. La journée simulée utilisée tout au long du chapitre compte 1\,555\,872 messages et 55\,542 transactions, soit 28 messages par
transaction~: la plupart des messages sont des cotations placées puis retirées sans jamais donner lieu à une transaction.

\begin{center}
\small
\begin{tabular}{@{}lp{4.4cm}p{7.2cm}@{}}
\toprule
type & champs & signification \\
\midrule
\texttt{A} ajout & heure, identifiant, côté, prix, quantité & un ordre à cours limité rejoint la fin de la file à son prix \\
\texttt{X} annulation & heure, identifiant, quantité & une partie ou la totalité d'un ordre au repos est retirée \\
\texttt{E} exécution & heure, identifiant, quantité, agresseur, identifiant de transaction & un ordre au repos est exécuté par un ordre entrant~; une ligne par ordre au repos exécuté \\
\bottomrule
\end{tabular}
\end{center}

Le mécanisme du simulateur est assez simple pour tenir en un paragraphe, et chaque chapitre ultérieur qui l'utilise s'appuie dessus. Un
prix efficient, que personne n'observe, saute d'un ou deux pas de cotation à des instants aléatoires. Les fournisseurs de liquidité ajoutent des ordres aux dix meilleurs
niveaux de chaque côté et les annulent à un rythme régulier, plus vite lorsque le prix efficient a traversé leur niveau. Les traders de bruit
envoient des ordres au marché dont les arrivées se regroupent (un processus de Hawkes, voir le livre~4, chapitre~7) et dont les signes persistent,
car la plupart sont des fragments d'ordres parents plus importants. Les traders informés envoient des ordres au marché en direction du prix efficient lorsque le
mid en est éloigné. Un niveau d'activité aléatoire, persistant sur plusieurs minutes, et un profil intrajournalier en U font varier ensemble tous ces
taux, et une rafale d'annonces en milieu de séance accélère le mouvement du prix efficient pendant quatre-vingt-dix secondes. Le simulateur
enregistre la vérité que les données réelles dissimulent~: le prix efficient, et quelles transactions étaient informées.

\begin{dated}{2026-09}{Quelle est la taille d'une journée de données ordre par ordre~?}\label{dat:rs:market-data-for-research:itch}
Nasdaq publie sur un serveur public des fichiers d'exemple de journées complètes de son flux TotalView-ITCH 5.0. Les fichiers compressés (gzip) pour
des journées isolées de 2019 et du début de 2020 pèsent de 3.5 à 5.6 milliards d'octets (30 janvier 2019~: 4.8~Go~; 30 janvier 2020~: 5.6~Go)~; les
fichiers publiés entre août 2025 et juin 2026 pèsent de 4.7 à 17.9~Go. Il s'agit du flux d'une seule bourse, pour chaque action qu'elle négocie, pour
une journée, compressé.
\end{dated}

\section{Horloges d'échantillonnage}

Une stratégie qui ne négocie pas à chaque message observe le marché à travers des \omterm{def:rs:market-data-for-research:bar}{barres}~: des synthèses des transactions comprises entre deux
instants. Ces instants sont fixés par une horloge, et celle-ci n'a pas besoin d'être l'horloge murale.

\begin{definition}[Barre, barre temporelle, barre de ticks, barre de volume, barre de montant, VWAP]\label{def:rs:market-data-for-research:bar}
Une \emph{barre}\index{barre} résume une suite de transactions consécutives par ses prix d'ouverture, plus haut, plus bas et de clôture, son volume, sa valeur
échangée et son nombre de transactions. Une \emph{barre temporelle}\index{barre temporelle} couvre un intervalle fixe de temps d'horloge murale~; une \emph{barre de
ticks}\index{barre de ticks} un nombre fixe de transactions~; une \emph{barre de volume}\index{barre de volume} se ferme à la transaction où son volume cumulé
atteint un seuil, et une \emph{barre de montant}\index{barre de montant} (dollar bar) se ferme à la transaction où sa valeur échangée cumulée atteint le sien.
Le \emph{prix moyen pondéré par les volumes}\index{prix moyen pondéré par les volumes} (VWAP) d'un ensemble de transactions est $\sum_i p_iq_i/\sum_i q_i$.
\end{definition}

\begin{definition}[Barre de déséquilibre]\label{def:rs:market-data-for-research:imbalance}
Une \emph{barre de déséquilibre}\index{barre de déséquilibre} se ferme lorsque la somme en valeur absolue des signes de transaction depuis son début, $|\theta| =
|\sum b_i|$, atteint un seuil fixé à partir des \omterm{def:rs:market-data-for-research:bar}{barres} déjà fermées~: le nombre espéré de transactions par \omterm{def:rs:market-data-for-research:bar}{barre} multiplié par
l'espérance du déséquilibre absolu par transaction (\omterm{def:rs:market-data-for-research:bar}{barre} de déséquilibre de ticks de López de Prado), minorée par la racine carrée du nombre espéré de transactions.
\end{definition}

\begin{omfigure}
\begin{tikzpicture}[font=\footnotesize,x=1.25cm,y=0.22cm]
\draw[->,>=Stealth] (0,0) -- (10.4,0) node[right] {temps};
\foreach \x/\h in {0.3/2,0.5/1,0.7/3,0.9/2,1.1/1,1.3/2,2.6/1,4.2/1,5.8/2,6.1/3,6.3/2,6.5/1,6.7/2,6.9/3,7.1/1,8.8/1,9.5/2}
  \draw[omDef,line width=1.6pt] (\x,0) -- (\x,\h*2);
\foreach \x in {2,4,6,8} \draw[black!60,dashed] (\x,-3.5) -- (\x,8.5);
\draw[omThm,very thick] (0.7,-8.5) -- (0.7,-4.5) (2.6,-8.5) -- (2.6,-4.5) (6.1,-8.5) -- (6.1,-4.5) (6.9,-8.5) -- (6.9,-4.5);
\node[black!70,anchor=west] at (0,10) {horloge temporelle~: une barre toutes les deux unités};
\node[omThm,anchor=west] at (0,-10.5) {horloge de volume~: une barre chaque fois que six unités de volume ont été échangées};
\node[omDef,anchor=west] at (7.2,10) {transactions (hauteur = quantité)};
\end{tikzpicture}

\omcaption{Les mêmes dix-sept transactions découpées par deux horloges. L'horloge temporelle ferme cinq intervalles égaux (en pointillés)~: trois presque
vides, l'une contient la rafale initiale, l'autre celle de la rafale entre $t = 6$ et 7. L'horloge de volume (en trait plein, en bas) ferme quatre
\omterm{def:rs:market-data-for-research:bar}{barres}~: deux courtes à l'intérieur des rafales, qui se terminent en 0.7 et 6.9, et deux longues qui absorbent les périodes calmes~; les trois
dernières transactions attendent la \omterm{def:rs:market-data-for-research:bar}{barre} suivante.}\label{fig:rs:market-data-for-research:schematic}
\end{omfigure}

Le plancher est nécessaire en pratique. Lorsque les achats et les ventes se compensent, le déséquilibre espéré par transaction est proche de zéro, et il en va de même du
seuil, et chaque transaction fermerait une \omterm{def:rs:market-data-for-research:bar}{barre}~; la racine carrée de $\E[T]$ est le niveau qu'une marche aléatoire de signes atteint pour la première fois après
environ $\E[T]$ transactions.

Pourquoi échantillonner selon l'activité plutôt que selon le temps~? Parce que la variance d'une variation de prix croît avec la quantité de transactions qui la sous-tend,
et que ces transactions sont inégalement réparties.

\begin{proposition}[Les barres temporelles sont un mélange de variances]\label{prop:rs:market-data-for-research:mixture}
Supposons que, étant donné le volume $V$ échangé dans une \omterm{def:rs:market-data-for-research:bar}{barre}, le rendement logarithmique de la \omterm{def:rs:market-data-for-research:bar}{barre} soit $\mathcal N(0, \sigma^2V)$. Alors le kurtosis des
rendements des \omterm{def:rs:market-data-for-research:bar}{barres temporelles} est
\[
\frac{\E[r^4]}{\E[r^2]^2} = \frac{3\,\E[V^2]}{\E[V]^2} = 3\bigl(1 + \mathrm{CV}(V)^2\bigr),
\]
avec $\mathrm{CV}$ le coefficient de variation des volumes des \omterm{def:rs:market-data-for-research:bar}{barres}, tandis que les \omterm{def:rs:market-data-for-research:bar}{barres} de même volume ont des rendements gaussiens.
\end{proposition}

\begin{proof}
$\E[r^2] = \sigma^2\E[V]$ et $\E[r^4] = 3\sigma^4\E[V^2]$ en conditionnant sur $V$. Une \omterm{def:rs:market-data-for-research:bar}{barre de volume} a un $V$ constant, à un dépassement près dû à la
dernière transaction.
\end{proof}

Clark (1973) a proposé cette subordination des prix au volume des transactions pour expliquer les rendements quotidiens à queue épaisse, et Ané et Geman
(2000) ont constaté que les rendements échantillonnés sur une horloge du nombre de transactions étaient proches de la loi gaussienne. Le jour simulé, les \omterm{def:rs:market-data-for-research:bar}{barres} d'une minute
dont les volumes ont un coefficient de variation de 0.94, de sorte que la proposition prédit un kurtosis de 5.67~; les rendements des mid d'une minute
ont un kurtosis de 5.44, et les rendements de clôture à clôture de 5.47. Les \omterm{def:rs:market-data-for-research:bar}{barres} découpées par les autres horloges, avec le même nombre de \omterm{def:rs:market-data-for-research:bar}{barres} dans la journée (environ
390), ont des kurtosis de 3.58 (tick), 4.04 (volume) et 3.92 (montant)~: pas 3, parce que dans le simulateur la volatilité suit le
niveau d'activité et non exactement le volume, mais on en est bien plus proche. Les horloges événementielles concentrent leurs \omterm{def:rs:market-data-for-research:bar}{barres} là où les transactions ont lieu
(la \cref{fig:rs:market-data-for-research:clocks})~: entre 2 et 42 \omterm{def:rs:market-data-for-research:bar}{barres} par quart d'heure, contre un rythme stable de 15 pour l'horloge temporelle.

\begin{omfigure}
\begin{tikzpicture}
\begin{axis}[width=11cm,height=5cm,xlabel={minutes après l'ouverture},ylabel={barres par quart d'heure},
  tick label style={font=\small},label style={font=\small},xmin=0,xmax=390,ymin=0,ymax=45,grid=major,grid style={gray!25},
  legend columns=3,legend style={font=\footnotesize,at={(0.5,-0.32)},anchor=north,draw=none,fill=none,
  /tikz/every even column/.append style={column sep=8pt}},table/col sep=comma]
\addplot[black!60,very thick,const plot mark mid] table[x=minute,y=time]{figdata/research/02-market-data-for-research/clocks.csv};
\addplot[omDef,very thick,const plot mark mid] table[x=minute,y=volume]{figdata/research/02-market-data-for-research/clocks.csv};
\addplot[omThm,thick,dashed,const plot mark mid] table[x=minute,y=tick]{figdata/research/02-market-data-for-research/clocks.csv};
\addplot[black,thin,sharp plot] coordinates {(210,0) (210,45)};
\legend{barres temporelles,barres de volume,barres de ticks}
\end{axis}
\end{tikzpicture}

\omcaption{\omterm{def:rs:market-data-for-research:bar}{Barres} fermées par quart d'heure sur la journée simulée, chaque horloge découpant environ 390 \omterm{def:rs:market-data-for-research:bar}{barres}. L'horloge temporelle ferme 15 \omterm{def:rs:market-data-for-research:bar}{barres}
à chaque quart d'heure~; les horloges de volume et de ticks suivent l'activité~: nourries après l'ouverture, calmes au milieu de la journée,
nourries autour de la rafale d'annonces à la minute 210 (ligne verticale) et avant la clôture. Les \omterm{def:rs:market-data-for-research:bar}{barres de montant} coïncident avec les \omterm{def:rs:market-data-for-research:bar}{barres de volume} à cette
échelle. Données~: \texttt{firm.tape}, la journée du chapitre, graine fixée.}\label{fig:rs:market-data-for-research:clocks}
\end{omfigure}

\section{Ce que chaque représentation détruit}

\begin{definition}[Série des mid]\label{def:rs:market-data-for-research:midquote}
Une \emph{série des mid}\index{série des mid} est le prix milieu $\tfrac12(b_t + a_t)$ échantillonné à des instants choisis, chaque valeur étant le
dernier mid à l'instant d'échantillonnage ou avant lui (un échantillon as-of).
\end{definition}

Un prix de transaction vaut le mid plus ou moins la moitié de la fourchette, selon l'initiateur de la transaction. Échantillonné trop finement, il rebondit.

\begin{proposition}[Le rebond ajoute une demi-fourchette au carré par rendement]\label{prop:rs:market-data-for-research:bounce}
Soit le prix de transaction $p = m + \tfrac s2 b$, avec une fourchette constante $s$ et des signes de transaction $b = \pm1$ indépendants entre eux et
du mid. Entre deux transactions, $\E[(\Delta p)^2] = \E[(\Delta m)^2] + s^2/2$.
\end{proposition}

\begin{proof}
$\Delta p = \Delta m + \tfrac s2(b_2 - b_1)$, et $\E[(b_2 - b_1)^2] = 2$ lorsque les signes sont indépendants et de moyenne nulle.
\end{proof}

C'est le modèle de Roll, dont l'estimateur de la fourchette est dérivé au livre~4, chapitre~21~; le signature plot de ce même chapitre trace la
variance réalisée en fonction de l'intervalle d'échantillonnage. Ici, le graphique répond à l'accroche
(la \cref{fig:rs:market-data-for-research:signature}). À une seconde, la variance du prix de transaction de la journée vaut 0.83 (en pour cent au
carré) et celle du mid 0.23~; à cinq minutes, elles valent 0.84 et 0.85. Le prix de transaction surestime à haute fréquence, comme l'annonce la
proposition. Le mid, lui, sous-estime, pour une raison que la proposition ne voit pas~: dans le simulateur, comme sur les marchés réels, le mid
rattrape le prix efficient par paliers, une cotation périmée à la fois, de sorte que ses rendements à court horizon sont positivement
autocorrélés et que leurs carrés s'additionnent pour donner moins que la variance sur des horizons plus longs. Aucune de ces deux séries n'est le prix efficient.

\begin{omfigure}
\begin{tikzpicture}
\begin{semilogxaxis}[width=10cm,height=5.2cm,xlabel={intervalle d'échantillonnage (secondes)},ylabel={variance réalisée (\%$^2$)},
  tick label style={font=\small},label style={font=\small},xmin=1,xmax=900,ymin=0,ymax=1.2,grid=major,grid style={gray!25},
  legend columns=2,legend style={font=\footnotesize,at={(0.5,-0.32)},anchor=north,draw=none,fill=none,
  /tikz/every even column/.append style={column sep=8pt}},table/col sep=comma]
\addplot[omThm,very thick,mark=*,mark size=1.2pt] table[x=seconds,y=trade]{figdata/research/02-market-data-for-research/signature.csv};
\addplot[omDef,very thick,mark=square*,mark size=1.2pt] table[x=seconds,y=mid]{figdata/research/02-market-data-for-research/signature.csv};
\legend{dernier prix de transaction,mid}
\end{semilogxaxis}
\end{tikzpicture}

\omcaption{Signature plot de la journée simulée~: variance réalisée des log-rendements en fonction de l'intervalle d'échantillonnage, à partir des prix de transaction
et à partir des mid. À une seconde, 0.83 contre 0.23~; à partir d'environ une minute, les deux s'accordent, et le bruit
de l'estimation sur une seule journée domine. Données~: \texttt{firm.tape}, la journée du chapitre, graine fixée.}\label{fig:rs:market-data-for-research:signature}
\end{omfigure}

Le rebond n'est qu'une perte parmi d'autres, et un entrepôt de recherche doit savoir laquelle chaque représentation subit~:
\begin{itemize}
\item \emph{Les prix de transaction} portent le rebond, et les transactions marquées de conditions particulières (lots irréguliers, comptes rendus hors séquence,
prix d'enchères, corrections tardives~: livre~1, chapitre~28) se mêlent aux transactions ordinaires. Le prix de l'enchère d'ouverture est un prix auquel un
volume important s'est échangé à un instant unique~; le traiter comme une transaction parmi d'autres fausse la première \omterm{def:rs:market-data-for-research:bar}{barre} de chaque journée.
\item \emph{Les mid} ne voient ni ce qui a été échangé ni en quelle quantité, et, lorsque la fourchette est large, le mid peut bouger sans aucune transaction.
\item \emph{Les \omterm{def:rs:market-data-for-research:bar}{barres}} perdent l'ordre des événements à l'intérieur de la \omterm{def:rs:market-data-for-research:bar}{barre}, le carnet d'ordres, et tous les événements entre deux échantillons~; le plus haut et le plus bas conservent une
trace du chemin suivi, que les estimateurs d'étendue du chapitre~7 exploitent.
\item \emph{Les \omterm{def:rs:market-data-for-research:bar}{barres temporelles} sur plusieurs instruments} sont alignées dans le temps mais pas dans l'information~: la dernière transaction d'un instrument peu liquide
peut dater de plusieurs minutes (transactions non synchrones et effet Epps, livre~4, chapitre~21).
\item \emph{Toute série horodatée par une seule horloge} masque l'écart entre l'horodatage de la bourse et l'instant où les données
sont arrivées~; un backtest doit utiliser le second (chapitre~3).
\end{itemize}

\section{Séries continues de contrats à terme}

Un contrat à terme arrive à échéance, et son successeur se négocie à un autre prix. Un historique assez long pour la recherche doit raccorder des
contrats successifs, et ce raccordement est un choix.

\begin{definition}[Série continue de contrats à terme, ajustement rétroactif, ajustement proportionnel]\label{def:rs:market-data-for-research:continuous}
Une \emph{série continue de contrats à terme}\index{série continue de contrats à terme} suit le contrat que détiendrait une position selon une règle de roll~:
le contrat le plus proche jusqu'à une date de roll quelques jours avant son dernier jour de négociation, puis le contrat suivant. L'\emph{ajustement rétroactif}\index{ajustement rétroactif} (back-adjustment)
ajoute, à chaque prix antérieur à un roll, l'écart entre le nouveau et l'ancien contrat à la date de roll, de sorte que la série n'ait aucun
saut au roll et que ses derniers prix sont réels. L'\emph{ajustement proportionnel}\index{ajustement proportionnel} (ratio adjustment) multiplie chaque prix antérieur par
le rapport entre le prix du nouveau contrat et celui de l'ancien à la date de roll, à la place.
\end{definition}

\begin{proposition}[Ce que préserve chaque ajustement]\label{prop:rs:market-data-for-research:adjust}
L'\omterm{def:rs:market-data-for-research:continuous}{ajustement proportionnel} préserve les rendements en pourcentage de la position roulée à chaque jour~; l'\omterm{def:rs:market-data-for-research:continuous}{ajustement rétroactif} préserve ses
variations de prix, le profit d'un contrat. Aucun des deux ne préserve les niveaux de prix passés, et les prix ajustés rétroactivement peuvent être négatifs.
\end{proposition}

\begin{proof}
Entre deux dates de roll, les deux ajustements agissent sur les prix du contrat détenu par une constante~: une constante additive laisse les différences
inchangées, une constante multiplicative laisse les rapports inchangés. Au passage d'une date de roll, l'ajustement est choisi de sorte que la variation de prix du nouveau contrat
(ou le rendement) soit enregistrée, et non le saut entre contrats. Les constantes additives s'accumulent avec les écarts de roll et peuvent dépasser
un prix passé.
\end{proof}

Les règlements-livraisons des contrats à terme sur pétrole brut publiés par la US Energy Information Administration rendent ce choix concret
(\cref{fig:rs:market-data-for-research:wti}). Du 2 janvier 2015 au 5 avril 2024, le contrat le plus proche est passé de USD~52.69 à
USD~86.91 le baril, en hausse de 65~\%. Une position longue, roulée cinq jours de négociation avant chacune des 111 échéances de la période, comme doit le faire un
fonds détenant le contrat le plus proche, a perdu 27~\%~: la plupart de ces années étaient en contango, et chaque roll vendait un contrat moins cher pour
acheter un contrat plus cher (le rendement de roll, dans le livre~3, chapitre~10). La série ajustée proportionnellement raconte l'histoire du fonds~; son premier prix est
USD~118.33. La série ajustée rétroactivement commence à USD~70.23 et, la backwardation de 2021--2023 ayant rendu négatifs les
écarts ultérieurs, atteint $-5.89$ le 21 avril 2020, jour où le contrat détenu réellement (juin 2020) a réglé à USD~11.57. La série
brute du contrat le plus proche affiche $-37.63$ le 20 avril 2020, soit le règlement d'un contrat arrivant à échéance qu'aucune position roulée ne détenait.

\begin{omfigure}
\begin{tikzpicture}
\begin{axis}[width=11cm,height=5.6cm,xlabel={année},ylabel={USD par baril},tick label style={font=\small},
  label style={font=\small},xmin=2015,xmax=2024.3,ymin=-40,ymax=130,grid=major,grid style={gray!25},
  xtick={2015,2016,2017,2018,2019,2020,2021,2022,2023,2024},xticklabels={2015,2016,2017,2018,2019,2020,2021,2022,2023,2024},
  scaled x ticks=false,legend columns=3,legend style={font=\footnotesize,at={(0.5,-0.3)},anchor=north,draw=none,fill=none,
  /tikz/every even column/.append style={column sep=8pt}},table/col sep=comma]
\addplot[black!45,thin] table[x=year,y=front]{figdata/research/02-market-data-for-research/wti.csv};
\addplot[omDef,thick] table[x=year,y=ratio]{figdata/research/02-market-data-for-research/wti.csv};
\addplot[omThm,thick,dashed] table[x=year,y=back]{figdata/research/02-market-data-for-research/wti.csv};
\legend{contrat le plus proche (tel que publié),ajusté proportionnellement,ajusté rétroactivement}
\end{axis}
\end{tikzpicture}

\omcaption{Contrats à terme sur pétrole brut WTI, 2015 à avril 2024~: le contrat le plus proche tel que publié (avec $-37.63$ le 20 avril 2020), et
série continue roulée cinq jours de négociation avant chaque échéance, ajustée proportionnellement (partant de 118.33) et ajustée rétroactivement (partant de
70.23, négative en avril 2020). Les trois séries se terminent au même prix. Données~: US Energy Information Administration, règlements-livraisons NYMEX,
contrats 1 et 2.}\label{fig:rs:market-data-for-research:wti}
\end{omfigure}

\begin{remark}[Quelle série pour quelle question]\label{rem:rs:market-data-for-research:which}
Rendements d'une stratégie roulée~: série ajustée proportionnellement. Profit d'un nombre fixe de contrats~: différences de la série ajustée rétroactivement. Un signal sur
niveaux de prix (une moyenne mobile du prix, une cassure au-dessus d'un plus haut passé)~: aucune des deux sans précaution, car les deux réécrivent le passé
à chaque roll~; calculer le signal sur le contrat détenu à ce moment-là. La forme de la courbe (carry, livre~8)~: les contrats non ajustés,
jamais une série ajustée.
\end{remark}

\section{Volumes, stockage et formats}

Une journée simulée d'un instrument de 1.56 million de messages à 36 octets (la taille d'un message ITCH add-order) représente 56~Mo avant
compression~; l'encadré daté montre ce que pèse une journée complète de bourse. Trois règles permettent de garder un entrepôt de recherche exploitable.

\begin{method}[Stocker les données de marché pour la recherche]\label{met:rs:market-data-for-research:store}
\begin{enumerate}
\item Conserver les messages bruts tels que reçus, immuables, avec leurs deux horodatages, partitionnés par date et par instrument~; chaque table
dérivée (carnets, \omterm{def:rs:market-data-for-research:bar}{barres}, variables) est reconstruite à partir d'eux par du code versionné.
\item Stocker les tables dérivées dans un format en colonnes, compressé, et les partitionner selon la manière dont elles sont lues~: par date pour
les analyses en coupe transversale, par instrument pour les analyses de séries temporelles.
\item Consigner pour chaque table dérivée la version du code, les partitions brutes qu'elle a lues et ses paramètres (l'horloge, le
seuil, la règle de roll), afin qu'un résultat puisse nommer les données exactes qu'il a utilisées (chapitre~29).
\end{enumerate}
\end{method}

\section{Tutoriel~: une journée, quatre horloges}\label{tut:rs:market-data-for-research:clocks}

\begin{tutorial}
\textbf{Objectif.} Simuler une journée de négociation, la découper selon quatre horloges et mesurer ce que chaque horloge et chaque prix font aux rendements.
\textbf{État final~:} les \cref{fig:rs:market-data-for-research:clocks,fig:rs:market-data-for-research:signature}~; kurtosis de 5.47 (temps),
3.58 (tick), 4.04 (volume), 3.92 (montant)~; le ratio de variance à une seconde, égal à 3.6.
\begin{tutsteps}
\item \textbf{La journée.} \texttt{simulate(TapeConfig(seconds=23\,400, u\_shape=1.5, news\_at=12\,600, seed=3))} renvoie les
messages, les transactions, le meilleur niveau du carnet après chaque message et le prix efficient (environ 10 secondes).
\item \textbf{\omterm{def:rs:market-data-for-research:bar}{Barres} à seuil.} Une \omterm{def:rs:market-data-for-research:bar}{barre} se clôt sur la transaction qui franchit le seuil~; les transactions ne sont jamais scindées entre deux \omterm{def:rs:market-data-for-research:bar}{barres}.
\omcode{code/firm/bars/firm_bars.py}{65}{90}{Les barres de volume et de montant se clôturent sur la transaction qui franchit le seuil.}\label{lst:rs:market-data-for-research:cuts}
\item \textbf{Deux prix, une journée.} Échantillonner la dernière transaction et le mid sur des grilles de 1 à 900 secondes et sommer les carrés
des rendements logarithmiques~; comparer le kurtosis à la prédiction du mélange issue des volumes des \omterm{def:rs:market-data-for-research:bar}{barres}.
\omcode{code/research/02-market-data-for-research/python/rs_marketdata.py}{45}{78}{Variance réalisée à partir des transactions et des mid~; le kurtosis du mélange.}\label{lst:rs:market-data-for-research:rv}
\item \textbf{Exécuter} \texttt{summary()} et \texttt{fig\_marketdata.py}.
\end{tutsteps}
\textbf{Pour aller plus loin.} Poser \texttt{act\_vol=0} (sans activité aléatoire) et observer le kurtosis des \omterm{def:rs:market-data-for-research:bar}{barres temporelles} décroître vers la
valeur du mélange de la seule forme en U~; découper des \omterm{def:rs:market-data-for-research:imbalance}{barres de déséquilibre} et suivre leur longueur au fil de la journée (exercice~7).
\end{tutorial}

\section{Construction~: le simulateur de marché et les constructeurs de barres}\label{bld:rs:market-data-for-research:tape}

\begin{build}
\bfield{Rôle} Les données intrajournalières standard de ce livre, avec la vérité jointe~: \texttt{firm.tape} écrit un flux ordre par ordre
à partir d'un mécanisme connu (prix efficient, flux informé et flux de bruit, annulations de cotations périmées)~; \texttt{firm.bars} transforme les transactions
en \omterm{def:rs:market-data-for-research:bar}{barres}. Les chapitres 8 à 10 mesurent des prédicteurs sur le tape, le chapitre~18 le rejoue, le chapitre~19 négocie contre lui.
\bfield{Interface} \texttt{TapeConfig} (durée, tick, niveaux, taux, paramètres de Hawkes et de métaordres, fréquence des sauts
du prix efficient, fenêtre d'annonce, profil intrajournalier, processus d'activité, graine)~; \texttt{simulate(cfg, v\_path)} renvoyant un \texttt{Tape} de
tableaux structurés \texttt{msgs}, \texttt{trades} (avec les vrais signes et un indicateur d'information), \texttt{top}, et les trajectoires \texttt{v\_t},
\texttt{v}, \texttt{act}~; \texttt{simulate\_pair(cfg, latency)}~; \texttt{Book().apply(msg)}. \texttt{time\_bars},
\texttt{tick\_bars}, \texttt{volume\_bars}, \texttt{dollar\_bars}, \texttt{imbalance\_bars}, \texttt{vwap}, \texttt{asof},
\texttt{continuous(c1, c2, expiry, days\_before)}.
\bfield{Règles} Déterministe pour une graine. Les taux sont constants par morceaux entre deux événements, sauf l'intensité de Hawkes, simulée exactement
par amincissement. Le carnet d'ouverture est un instantané à l'instant zéro. Les prix sont exprimés en nombres entiers de ticks, les tailles en nombres entiers d'actions.
\bfield{Tests d'acceptation} \texttt{code/firm/tape/tests/} et \texttt{code/firm/bars/tests/}~: sortie identique pour une graine~;
les messages reconstruisent le meilleur niveau enregistré du carnet à chaque étape et le carnet ne se croise jamais~; les exécutions égalent les transactions~; le mid suit
le prix efficient~; agrégé sur six graines, le déséquilibre prédit le mouvement du mid suivant et les transactions informées font perdre de l'argent au
côté passif, tandis que les transactions de bruit le paient~; les \omterm{def:rs:market-data-for-research:bar}{barres} se clôturent sur la transaction de franchissement~; un roll vérifié à la main.
\bfield{Pour aller plus loin} Plusieurs instruments partageant un facteur dans leurs prix efficients~; des enchères à l'ouverture et à la clôture~; un portage en C++20
de la boucle d'événements.
\end{build}

\begin{omsources}
\item P.~K.~Clark, ``A subordinated stochastic process model with finite variance for speculative prices'',
\emph{Econometrica} 41(1), 1973.
\item T.~Ané et H.~Geman, ``Order flow, transaction clock, and normality of asset returns'', \emph{Journal of Finance} 55(5),
2000.
\item D.~Easley, M.~López de Prado et M.~O'Hara, ``The volume clock: insights into the high-frequency paradigm'',
\emph{Journal of Portfolio Management} 39(1), 2012.
\item M.~López de Prado, \emph{Advances in Financial Machine Learning}, Wiley, 2018, chapitre 2 (barres guidées par l'information).
\item R.~Roll, ``A simple implicit measure of the effective bid--ask spread in an efficient market'', \emph{Journal of Finance}
39(4), 1984.
\item Nasdaq, spécification de \emph{TotalView-ITCH 5.0}, et les fichiers d'exemple sur \texttt{emi.nasdaq.com}.
\item US Energy Information Administration, prix des contrats à terme NYMEX, contrats sur pétrole brut 1--4 (quotidiens).
\end{omsources}

\section{Exercices}

\begin{exercise}[$\star$]\label{exo:rs:market-data-for-research:1}
Transactions de 300 actions à 20.00, 500 à 20.02 et 200 à 19.99. Quel est leur VWAP~?
\end{exercise}

\begin{exercise}[$\star$]\label{exo:rs:market-data-for-research:2}
Une action cotée à USD~20 présente une fourchette d'un cent, se négocie 20\,000 fois par jour et affiche une volatilité quotidienne vraie de 2~\%. D'après
la \cref{prop:rs:market-data-for-research:bounce}, quelle volatilité quotidienne la somme des carrés des rendements logarithmiques transaction à transaction affiche-t-elle~?
\end{exercise}

\begin{exercise}[$\star$]\label{exo:rs:market-data-for-research:3}
Un flux représente en moyenne 40 octets par message, et un instrument actif émet 1.5 million de messages par jour. Quelle quantité brute de données produisent
5\,000 instruments de ce type sur une année de 252 jours~?
\end{exercise}

\begin{exercise}[$\star\star$]\label{exo:rs:market-data-for-research:4}
Les volumes à la minute ont un coefficient de variation de 0.8. Quelle kurtosis prédit la \cref{prop:rs:market-data-for-research:mixture}
pour les rendements à la minute~? Quel coefficient de variation donne une kurtosis de 6~?
\end{exercise}

\begin{exercise}[$\star\star$]\label{exo:rs:market-data-for-research:5}
Le contrat de l'échéance la plus proche se règle à 70, 71 et 72 sur trois jours, et le suivant à 73, 74 et 76~; la position est transférée lors du roll à la clôture du
deuxième jour. Donner la série détenue, puis celle issue de l'\omterm{def:rs:market-data-for-research:continuous}{ajustement rétroactif} (back-adjustment) et celle de l'\omterm{def:rs:market-data-for-research:continuous}{ajustement proportionnel} (ratio adjustment).
\end{exercise}

\begin{exercise}[$\star\star$]\label{exo:rs:market-data-for-research:6}
Une action se négociait à USD~10 en 2015 et à USD~100 en 2025, pour une même valeur échangée chaque jour. Un seuil de \omterm{def:rs:market-data-for-research:bar}{barres de volume} exprimé en
actions donne combien de \omterm{def:rs:market-data-for-research:bar}{barres} par jour en 2015 par rapport à 2025~? Et pour un seuil de \omterm{def:rs:market-data-for-research:bar}{barres de montant} (dollar bars)~?
\end{exercise}

\begin{exercise}[$\star\star\star$]\label{exo:rs:market-data-for-research:7}
\emph{Programmation.} Découper la journée du chapitre en \omterm{def:rs:market-data-for-research:imbalance}{barres de déséquilibre}, à partir d'un nombre attendu de 142 transactions par \omterm{def:rs:market-data-for-research:bar}{barre}. Combien de \omterm{def:rs:market-data-for-research:bar}{barres} se ferment, combien
de transactions contiennent-elles en moyenne, et quelle est la kurtosis de leurs rendements des mid~? Pourquoi les \omterm{def:rs:market-data-for-research:bar}{barres} se resserrent-elles au fil de la journée~?
\end{exercise}

\begin{exercise}[$\star\star\star$]\label{exo:rs:market-data-for-research:8}
\emph{Trouver la faille.} «~Notre signal de carry est le rapport logarithmique annualisé entre la série du deuxième mois à \omterm{def:rs:market-data-for-research:continuous}{ajustement rétroactif} et
la série de l'échéance la plus proche à \omterm{def:rs:market-data-for-research:continuous}{ajustement rétroactif}, et notre signal de tendance est la variation en pourcentage de cette dernière sur
douze mois.~»
\end{exercise}

\section{Problème~: Quelle horloge~?}

\begin{problem}[{Problème du week-end --- une journée, quatre horloges et deux prix}]\label{pb:rs:market-data-for-research:1}
La journée simulée du chapitre~: 6.5 heures, un profil en U, un niveau d'activité aléatoire et une rafale de nouvelles à 12\,600 secondes, graine 3.

\medskip\noindent\textbf{Partie I --- Le flux.}
\begin{enumerate}
\item Combien de messages et de transactions compte la journée, et combien de messages par transaction~?
\item Quelle est, à peu près, la part des ajouts, des annulations et des exécutions parmi les messages, et pourquoi les ajouts
et les annulations sont-ils presque égaux~?
\item Quelle est la taille brute de la journée à 36 octets par message~?
\item Pourquoi un entrepôt de recherche conserve-t-il les messages plutôt que les \omterm{def:rs:market-data-for-research:bar}{barres}~?
\end{enumerate}

\medskip\noindent\textbf{Partie II --- Horloges.}
\begin{enumerate}[resume]
\item Combien de \omterm{def:rs:market-data-for-research:bar}{barres} chaque horloge découpe-t-elle, avec des seuils choisis pour environ 390~?
\item Quelles sont les kurtosis des rendements clôture à clôture pour les quatre horloges~?
\item Quel est le coefficient de variation des volumes des \omterm{def:rs:market-data-for-research:bar}{barres} d'une minute, et quelle kurtosis la proposition relative au mélange
prédit-elle pour elles~?
\item Pourquoi les kurtosis des horloges événementielles sont-elles supérieures à 3~?
\item Où dans la journée se concentrent les \omterm{def:rs:market-data-for-research:bar}{barres de volume}, et pourquoi~?
\end{enumerate}

\medskip\noindent\textbf{Partie III --- Prix.}
\begin{enumerate}[resume]
\item Quelles sont les variances réalisées à partir des transactions et des mid, à 1, 5, 60 et 300 secondes~?
\item Quel effet rend la série des transactions trop élevée à une seconde~?
\item Lequel rend la \omterm{def:rs:market-data-for-research:midquote}{série des mid} trop basse~?
\item À partir de quel intervalle les deux séries concordent-elles à moins de 5~\%~?
\item Quel estimateur employer pour évaluer la variance de la journée à partir des données à la seconde (livre~4, chapitre~21)~?
\end{enumerate}

\medskip\noindent\textbf{Partie IV --- Jugement.}
\begin{enumerate}[resume]
\item Quelle horloge retenir pour une estimation de la volatilité quotidienne, et laquelle pour un modèle des mouvements de prix à court terme~?
\item Qu'est-ce qu'une \omterm{def:rs:market-data-for-research:bar}{barre temporelle} d'une minute perd et qu'une \omterm{def:rs:market-data-for-research:bar}{barre de volume} conserve~?
\item Qu'est-ce qu'une \omterm{def:rs:market-data-for-research:bar}{barre de volume} perd et qu'une \omterm{def:rs:market-data-for-research:bar}{barre temporelle} conserve~?
\item Énoncer le \emph{résultat nommé}~: la kurtosis des rendements des \omterm{def:rs:market-data-for-research:bar}{barres temporelles} d'une minute, comparée à celle des rendements des \omterm{def:rs:market-data-for-research:bar}{barres de montant} sur la journée, et le
rapport de variance transactions/mid à une seconde~?
\item Qu'est-ce qui, dans le simulateur, rend les \omterm{def:rs:market-data-for-research:bar}{barres temporelles} à queue épaisse, et qu'est-ce qui la supprimerait~?
\item En une phrase~: quelle représentation est «~le~» prix~?
\end{enumerate}
\end{problem}

\section{Questions d'entretien}

\begin{interviewq}[$\star$ \iqroles{chercheur, trader}]\label{iq:rs:market-data-for-research:1}
Pourquoi pourriez-vous échantillonner les prix selon le volume plutôt que selon le temps~?
\end{interviewq}

\begin{interviewq}[$\star\star$ \iqroles{chercheur}]\label{iq:rs:market-data-for-research:2}
Votre volatilité réalisée calculée sur les prix des transactions à la seconde est le double de celle calculée sur les prix à cinq minutes. Que se passe-t-il~?
\end{interviewq}

\begin{interviewq}[$\star\star$ \iqroles{chercheur, trader}]\label{iq:rs:market-data-for-research:3}
Comment construiriez-vous une \omterm{def:rs:market-data-for-research:continuous}{série continue de contrats à terme} pour un backtest de suivi de tendance, et qu'est-ce qui ne va pas avec les prix à \omterm{def:rs:market-data-for-research:continuous}{ajustement rétroactif}~?
\end{interviewq}

\begin{interviewq}[$\star\star$ \iqroles{développeur, chercheur}]\label{iq:rs:market-data-for-research:4}
Concevez le stockage d'une année de données ordre par ordre pour 5\,000 instruments, à usage de recherche.
\end{interviewq}

\begin{interviewq}[$\star\star$ \iqroles{chercheur, ingénieur ML}]\label{iq:rs:market-data-for-research:5}
Le modèle d'un collègue utilise des \omterm{def:rs:market-data-for-research:bar}{barres} d'une minute sur cinquante actions. Qu'est-ce qui peut aller de travers parce que ces \omterm{def:rs:market-data-for-research:bar}{barres} sont temporelles~?
\end{interviewq}

\begin{interviewq}[$\star\star\star$ \iqroles{chercheur}]\label{iq:rs:market-data-for-research:6}
Montrez que si les rendements sont gaussiens conditionnellement au volume échangé, les rendements des \omterm{def:rs:market-data-for-research:bar}{barres temporelles} ont une kurtosis de $3(1 + \mathrm{CV}^2)$.
\end{interviewq}
